\section{Pakes, the Lambert law and the Thorin condition}\label{sec:pakes}

Anthony Pakes has worked on the objects that surround the Thorin condition: size biasing and the classical moment problem, the Laplace transforms of positive infinitely divisible laws, the Lambert function, and ladders of self decomposable laws. This section follows the first author's account of that work and of what it gives for the clock \cite{P26pak}, verified, corrected where needed, and written in the notation of the book. In Pakes's language the question of Section~\ref{ch:ggc} is whether $\log G$ is a Thorin Bernstein function, that is, a Bernstein function whose L\'evy measure has the form $k(t)\,dt/t$ with $k$ completely monotone \cite[Chapter 8]{boSSV}, or equivalently the Laplace exponent of a GGC. By Theorem~\ref{bh:thm:main} this holds if and only if $(\log G)'=G'/G$ is a Stieltjes function, and if and only if RH holds. For the Lambert function $\mathcal W$ the corresponding statement, that $\mathcal W'$ is a Stieltjes function, is a theorem, and the Thorin measure of $\mathcal W$ is explicit. The section follows this parallel as far as it goes.

Section~\ref{pk:sec:strands} reviews the two strands of Pakes's work. Section~\ref{pk:sec:form} writes the Thorin question in his form: with no hypothesis the size biased clock is the clock plus an independent variable whose density is proportional to the heat trace, and RH holds exactly when this variable is a mixture of exponential laws; the moment problem of his first strand becomes the Hankel test of Section~\ref{ch:heat}, whose negative eigenvalues count the zeros off the line (Theorem~\ref{thm:ch12:inertia}), and Table~\ref{tab:ch21:nfaillong} follows the count along the de Bruijn Newman flow. Section~\ref{pk:sec:lambert} proves the theorem of Kalugin, Jeffrey and Corless that the Lambert law is a GGC with an explicit Thorin measure, by boundary values on the cut, reflection and a growth bound, and isolates the one step of that proof which does not transfer to the clock. The first author's account then localises, with no hypothesis, the points of the upper half plane where $\im(G'/G)$ could fail to be negative, as it must be for a Stieltjes function, to a thin parabola along the negative axis beyond about $-T_0^2$, where $T_0=3\times10^{12}$ is the height to which RH is verified; this is Lemma~\ref{bs:lem:disc} and Theorem~\ref{bs:thm:local} in Section~\ref{ch:boschB}, and what the zeros off the line could contribute there is discussed in Section~\ref{ag:sec:exactly}. Section~\ref{pk:sec:ladder} places the clock among Urbanik's multiply self decomposable laws, in $L_2$ and not in $L_3$ with no hypothesis, and shows that it is not the martingale limit of a supercritical continuous state branching process with finite Malthusian parameter. Section~\ref{pk:sec:filter} derives from the theorem of Rodgers and Tao a filter for the properties of the P\'olya kernel that could imply RH, and Section~\ref{pk:sec:assess} sorts what survives and lists open problems.

We do not prove RH, and several of the equivalences are elementary. What the section adds is a dictionary between Pakes's language and that of the book, a few unconditional facts about the clock, the solved Lambert problem as a template whose single failing step for the clock is RH itself, a list of routes that are closed, and one property, the GGC property of the variable of Biane, Pitman and Yor, that passes the filter of Section~\ref{pk:sec:filter} wherever it has been tested, but fails for most of the real Dirichlet $L$ functions tested there, so that it does not detect the Euler product. The programs \texttt{pakes\_sizebias.py} and \texttt{pakes\_lambert.py} compute the numbers of Sections~\ref{pk:sec:form} and \ref{pk:sec:lambert}.

\subsection{Two strands of Pakes's work}\label{pk:sec:strands}

The first strand is the classical moment problem. For a random variable $X\ge0$ with $0<\E X<\infty$, the size biased variable $\widehat X$, length biased in Pakes's word, has the law $x\,P(X\in dx)/\E X$. Pakes and Khattree \cite{PK92} characterise laws by in law equations linking $\widehat X$ and $X$, such as $X\stackrel d=\mathsf U\widehat X$ with $\mathsf U$ uniform on $(0,1)$ and independent of $\widehat X$, which holds exactly for the exponential laws, and they study the role of the moment problem in such characterisations, which can single out laws that are not determined by their moments. Pakes \cite{Pakes96} studies size biasing for the laws with the moments of the log normal law, and in \cite{Pakes01} he proves converse forms of the Carleman and Krein criteria and introduces size biasing and the stationary excess operation of order $r$; Berg \cite{Berg04} showed, correcting a remark in \cite{Pakes01}, that a law can be determinate in the sense of Stieltjes while its size biased and stationary excess versions of order one are not. Pakes, Hung and Wu \cite{PHW01} give criteria for the determination of a law by its moments, and Pakes \cite{Pakes07} describes the structure of the Stieltjes classes of moment equivalent laws; these papers are cited throughout the later literature on the moment problem, for instance by Stoyanov and his coauthors and by Penson and his coauthors.

Size biasing meets infinite divisibility in Steutel's theorem: $X$ is infinitely divisible if and only if $\widehat X\stackrel d=X+A$ for some $A\ge0$ independent of $X$, and the law of $A$ is then unique \cite{Steutel73}, \cite[Theorem 51.1]{Sato99}, \cite{SvH}; in this form the theorem is \cite[Theorem 11.1]{AGK19}, and Pakes takes it for positive laws from Pakes, Sapatinas and Fosam \cite{PSF96}. In Laplace transforms, with $\Psi(u)=-\log\E e^{-uX}$, it reads $\E e^{-u\widehat X}=\E e^{-uX}\,\Psi'(u)/\Psi'(0)$, so that $A$ has the Laplace transform $\Psi'(u)/\Psi'(0)$. In \cite{Pakes14} Pakes assumes the decomposition and asks when $A$ or $\widehat X$ is again infinitely divisible; he shows that size biasing up to the fourth order preserves infinite divisibility for the Lambert law and for a compound Poisson law, and conjectures that it does so at every order. Section~\ref{pk:sec:form} applies Steutel's theorem to the clock.

For the laws of this book the moment problem is determinate. The moment generating function $\xi(\frac12+z)/\xi(\frac12)$ of the P\'olya density is entire \eqref{eq:ch2:mgf}. The clock has $\E e^{sT}<\infty$ exactly for $s<\gamma_1^2=199.7904548\ldots$: every zero has $|\rho-\frac12|\ge|\im\rho|\ge\gamma_1$, so $1/G$ is analytic in the disc $|u|<\gamma_1^2$ and has a pole at $-\gamma_1^2$, and by Landau's theorem \cite[Chapter~II, Theorem~5b]{Widder41} the Laplace transform of a positive measure is singular at the real point of its abscissa of convergence, which is therefore $-\gamma_1^2$. Under RH the measure $\nu=\sum_ka_k^{-1}\delta_{1/a_k}$ of the Stieltjes moments $m_n$ is supported in $[0,\gamma_1^{-2}]$, so its moment problem is a determinate Hausdorff problem (Section~\ref{ch:heat}), while the Thorin measure $U_{\RH}$ itself has unbounded support and infinite mass.

The second strand concerns the Laplace transforms of positive infinitely divisible laws, and the Lambert function $\mathcal W$, the principal solution of $\mathcal W(u)e^{\mathcal W(u)}=u$ (Section~\ref{pk:sec:lambert}). Pakes \cite{Pakes11} and, independently, Kalugin and Jeffrey \cite{KJ10} proved that $\mathcal W$ is a Bernstein function, from a formula for the derivatives of $\mathcal W$ that involves a sequence of polynomials \cite[(8)]{Pakes26}: Pakes through the positivity of the polynomials, Kalugin and Jeffrey by showing that their coefficients form positive, log concave and unimodal sequences. So $e^{-\mathcal W(u)}=\mathcal W(u)/u$ is the Laplace transform of an infinitely divisible law, the Lambert law. In the same paper, by Pakes's own account \cite[Section 2]{Pakes26}, a complex variable argument shows that $\mathcal W'$ is a Stieltjes function, so that $\mathcal W$ is a Thorin Bernstein function and the Lambert law is a GGC. Kalugin, Jeffrey, Corless and Borwein \cite{KJCB12} obtained Stieltjes representations independently and showed that $\mathcal W$ is a complete Bernstein function; the distribution function of the Thorin measure of $\mathcal W$ is $\pi^{-1}\im\mathcal W(-y+i0)$, which vanishes for $y<1/e$ \cite{KJCB12}, \cite[Theorem 3.3]{KJC11}, and the Thorin representation itself is \cite[Theorem 6.2]{KJC11}, published, with integral representations of other types, in \cite{KJC12}. Pakes gave two further proofs that $\mathcal W$ is Thorin Bernstein: from Nuttall's integral for the coefficients $n^n/n!$ of $\mathcal W'$, which yields an explicit quantile function of the Thorin measure \cite{Pakes18}, and from continuous state branching processes, one of whose martingale limits has the cumulant function $\mathcal W$ \cite{Pakes20}, which also gives the self decomposability of the Lambert law.

In \cite{Pakes26} Pakes identifies the background driving L\'evy process of the Lambert law: it is compound Poisson with unit rate and jump density $\int ze^{-xz}\,U_{\mathcal W}(dz)$, where $U_{\mathcal W}$ is the Thorin measure of $\mathcal W$. With $T_{\mathcal W}$ a variable with the Lambert law and $J$ a jump of the driving process, this gives the in law relations $T_{\mathcal W}\stackrel d=\mathsf U(T_{\mathcal W}+J)$ and $T_{\mathcal W}\stackrel d=\mathsf U\,\mathcal SJ$, with $\mathsf U$ uniform on $(0,1)$, $\mathcal S$ the stationary excess operator and the factors on the right independent; and if both relations hold for a variable $X$ and some $J\ge0$ of finite mean $\mu>0$, then $X\stackrel d=\mu T_{\mathcal W}$, so that the two relations characterise the Lambert law up to a scale factor \cite[Theorem 1]{Pakes26} (Remark~\ref{pk:rem:sd}). He embeds the Lambert law in a family of GGCs whose Laplace transforms solve Lambert's trinomial equation \cite[Theorem 6]{Pakes26}, which gives one more proof that $\mathcal W$ is Thorin Bernstein \cite[Corollary 2]{Pakes26}, and he introduces BDLP ladders: from an infinitely divisible $X_0\ge0$ with $\E(\log(1+X_0))^n<\infty$ for every $n$, the law of $X_{n+1}$ is the self decomposable law driven by the L\'evy process of $X_n$ \cite[Section 5]{Pakes26}. In his indexing $X_1$ is self decomposable; this is a unit shift from Urbanik's classes $L_n$ used in Section~\ref{pk:sec:ladder}, in which $L_0$ is the class of self decomposable laws \cite{Urbanik73,Sato80,BergForst83}. Along a ladder he finds the size biased decomposition $\widehat X_{n+1}\stackrel d=X_{n+1}+A_n$, with $A_n$ of density proportional to the tail of the L\'evy measure of $X_n$ \cite[(53)]{Pakes26}, and this is the form in which Section~\ref{pk:sec:form} applies Steutel's theorem to the clock.

\subsection{The Thorin question in Pakes's form}\label{pk:sec:form}

For the clock the objects of both strands are at hand. The Thorin question asks whether $\log G$ is a Thorin Bernstein function, that is, whether $(\log G)'=G'/G$ is a Stieltjes function (Theorem~\ref{bh:thm:main}); then the Thorin measure is $U_{\RH}=\sum_{\gamma>0}\delta_{\gamma^2}$, with no drift (Proposition~\ref{prop:rigid}), and Stieltjes inversion \eqref{eq:ch3:inversion} recovers it from the boundary values of $G'/G$ on the cut, or, in integrated form, from those of $\log G$ (Section~\ref{sec:ch3:readings}). With no hypothesis, $\log G$ is a Bernstein function with L\'evy measure $W(t)\,dt/t$, where $W(t)=\sum_ke^{-a_kt}$ is the heat trace (Theorem~\ref{thm:ch12:idsd}); under RH, $W(t)=\int e^{-tz}\,U_{\RH}(dz)$, the counterpart for the clock of the formula that gives the canonical function of the Lambert law as the Laplace transform of the density of its Thorin measure (Corollary~\ref{pk:cor:levy}; \cite[(12)]{Pakes26}).

The Bernstein property is weaker than the Thorin Bernstein property, but it is not blind to the zeros off the line. A Bernstein function extends holomorphically to $\re u>0$, and the zeros $-a$ of $G$ that come from a zero $\rho=\frac12+\delta+i\gamma$ off the line have $\re(-a)=\delta^2-\gamma^2<0$, so holomorphy alone cannot see them. But the Bernstein property of $\log G$ says more, $W(t)\ge0$ for every $t>0$, and that would exclude, for instance, a zero off the line below all the other zeros: if $\re a=\gamma^2-\delta^2$ were smaller than $\re a_k$ for every other $a_k$, which is $\gamma_k^2$ for a zero on the line, the four zeros $\rho$, $1-\rho$, $\bar\rho$, $1-\bar\rho$ would contribute $2e^{-(\gamma^2-\delta^2)t}\cos(2\delta\gamma t)$ to $W(t)$, which dominates as $t\to\infty$ and changes sign. The property holds with no hypothesis (Theorem~\ref{thm:ch12:idsd}, whose proof uses the verification of RH below height $1000$), so it cannot decide RH; among the Bernstein functions it is the Thorin Bernstein property that is equivalent to RH.

Size biasing gives the Thorin question a form with one variable. The size biased clock $\widehat T$ has the law $t\,P(T\in dt)/\E T$, and $\E T=m_0=\sum_ka_k^{-1}=0.0231049931\ldots$, the first Stieltjes moment of Section~\ref{ch:heat}.

\begin{proposition}\label{pk:prop:sizebias}
With no hypothesis, $\widehat T\stackrel d=T+A$, with $A$ independent of $T$ and of density $W(t)/m_0$ on $(0,\infty)$. RH holds if and only if $A$ is a mixture of exponential laws, and then $A\stackrel d=E/V$, with $E$ standard exponential and independent of $V$, and $V$ of law $m_0^{-1}\sum_{\gamma>0}\gamma^{-2}\delta_{\gamma^2}$, the ordinates counted with multiplicity.
\end{proposition}

\begin{proof}
The Laplace transform of $\widehat T$ is
\[
\frac{\E[Te^{-uT}]}{m_0}=-\frac1{m_0}\,\frac d{du}\,\frac1{G(u)}=\frac1{m_0}\,\frac1{G(u)}\,\frac{G'}G(u).
\]
By Section~\ref{ch:clock}, $(G'/G)(u)=\int_0^\infty e^{-ut}W(t)\,dt$ for $u>0$ with no hypothesis, and $W>0$ (Theorem~\ref{thm:ch12:idsd}); letting $u\downarrow0$, monotone convergence gives $\int_0^\infty W(t)\,dt=(G'/G)(0)=G'(0)=\sum_ka_k^{-1}=m_0$. So $W/m_0$ is a probability density with Laplace transform $(G'/G)(u)/m_0$, the Laplace transform of $\widehat T$ is the product of those of $T$ and of a variable $A$ with this density, and by the uniqueness of Laplace transforms $\widehat T\stackrel d=T+A$ with $A$ independent of $T$. A variable $A>0$ is a mixture of exponential laws, $A\stackrel d=E/V'$ with $V'>0$ independent of $E$, exactly when it has a completely monotone density: the density of $E/V'$ is $\E[V'e^{-V't}]$, and conversely, by Bernstein's theorem, a completely monotone probability density is $\int_{[0,\infty)}e^{-vt}\,\sigma(dv)$ with $\int v^{-1}\sigma(dv)=1$, which forces $\sigma(\{0\})=0$, and this is the density of $E/V'$ with $P(V'\in dv)=v^{-1}\sigma(dv)$. The density $W/m_0$ is completely monotone exactly when RH holds (Section~\ref{ch:clock} and Table~\ref{tab:ch3:hierarchy}), and then $W(t)=\sum_{\gamma>0}e^{-\gamma^2t}=m_0\E[Ve^{-Vt}]$ with $V$ as stated.
\end{proof}

This is Steutel's decomposition of Section~\ref{pk:sec:strands} with $\Psi=\log G$, and it is the identity that Pakes derives along BDLP ladders \cite[(53)]{Pakes26}, here for the clock and the subordinator of Proposition~\ref{sb:prop:bdlp} that drives it, whose L\'evy density is $h=-W'$. Pakes's variable $A_n$ has the density $x\mapsto\int_x^\infty\omega_n(y)\,dy$, the tail of the L\'evy density $\omega_n$ of the driving law, normalised by the mean $\int_0^\infty x\,\omega_n(x)\,dx$, and for the clock the tail is $\int_t^\infty h(x)\,dx=W(t)$ and the mean is $\int_0^\infty xh(x)\,dx=\int_0^\infty W(t)\,dt=m_0$. The content of Proposition~\ref{pk:prop:sizebias}, RH if and only if $W$ is completely monotone, is that of Proposition~\ref{sb:prop:bdlp}; the size biased form puts it in one variable that exists with no hypothesis. Its mean is $\E A=m_1/m_0=\mathrm{Var}\,T/\E T=1.60885567\times10^{-3}$, from the Taylor coefficients of $\xi$ at $\frac12$, with no zero entering. Under RH, $A$ is infinitely divisible, as a mixture of exponential laws \cite{Goldie67,Steutel67}, and so is $\widehat T=T+A$; so under RH both parts of the question that Pakes asks in \cite{Pakes14} have a positive answer for the clock. The weights $P(V=\gamma^2)=\gamma^{-2}/m_0$ of the first six ordinates are $0.21663$, $0.09794$, $0.06919$, $0.04676$, $0.03990$ and $0.03064$: the lowest zero carries about a fifth of the law of $V$, and the first six carry $0.50105$, about a half. The first $1700$ ordinates carry $0.978508$. The rest is $m_0^{-1}\int_{(\gamma_{1700},\infty)}t^{-2}\,dN(t)$; the smooth part $\vartheta(t)/\pi+1$ of the count, with the boundary term $(\vartheta(\gamma_{1700})/\pi+1-1700)/(m_0\gamma_{1700}^2)$ that it leaves at $\gamma_{1700}$, accounts for it to within $3\times10^{-9}$, and the remainder is $2m_0^{-1}\int_{\gamma_{1700}}^\infty S(t)\,t^{-3}\,dt$, where $S(t)$ is the fluctuating part of the count. The program \texttt{pakes\_sizebias.py} computes these numbers, the moments at $80$ digits and the weights from the $70$ digit ordinates of \texttt{zeros\_1700.txt}, checked against mpmath's \texttt{zetazero}.

Pakes's first strand, the moment problem, enters through the Taylor coefficients at the origin. For $|u|<\gamma_1^2$ and for $|u|<1/e$ respectively,
\[
(\log G)'(u)=\sum_{n\ge0}m_n(-u)^n,\qquad \mathcal W'(u)=\sum_{n\ge1}\frac{n^n}{n!}(-u)^{n-1},
\]
with the Stieltjes moments $m_n=\sum_ka_k^{-n-1}$ of Section~\ref{ch:heat}; the second series is the derivative of the Lagrange series of $\mathcal W$ \cite[(3.1)]{Corless96}. Under RH, $m_{n-1}=\int z^{-n}\,U_{\RH}(dz)$ for $n\ge1$, just as $n^n/n!=\int z^{-n}\,U_{\mathcal W}(dz)$ (Corollary~\ref{pk:cor:nuttall}): both coefficient sequences are then moment sequences of positive measures, the images of $z^{-1}U_{\RH}(dz)$ and of $U_{\mathcal W}$ under $z\mapsto1/z$, and their Hankel matrices are positive definite. For $\mathcal W$ this is a consequence of the theorem of Section~\ref{pk:sec:lambert}. For the clock the positive definiteness of every $\mathcal H_N=[m_{i+j}]_{i,j<N}$ is equivalent to RH (Section~\ref{ch:heat}); the number of negative eigenvalues of $\mathcal H_N$ does not decrease with $N$ and tends to the number of distinct zeros $\rho$ with $\re\rho>\frac12$ and $\im\rho>0$ (Theorem~\ref{thm:ch12:inertia}); and Table~\ref{tab:ch21:nfaillong} shows the count at work along the de Bruijn Newman flow, where in the data a pair of zeros that starts at $(\gamma_k,\gamma_{k+1})$ is often registered soon after its collision at an order close to $2k+2$, although this rule fails at order $50$ (Section~\ref{sec:ch21:hankel}).

Li's criterion is a positivity condition on the same zeros, read on the Thorin functional through the test functions $\varphi_n$ of Section~\ref{sec:ch20:thorin}. It sees a zero off the line more slowly: the term of $\rho=\frac12-d+i\gamma$, with $0<d<\frac12$, grows like $e^{nd/\gamma^2}$, so it can dominate only from $n$ of order $\gamma^2/d$, which exceeds $2\gamma^2$, while in the data of Section~\ref{sec:ch21:hankel} the Hankel test often registers a pair at an order close to $2k+2$, set by its index $k$, at a cost in precision that grows linearly with the order (Section~\ref{ch:heat}). Both are tests at low height, and in both every zero off the line is eventually seen.

\subsection{The Lambert law, a Thorin problem solved}\label{pk:sec:lambert}

Let $\mathcal W$ be the principal branch of Lambert's function, the solution of $\mathcal W(u)e^{\mathcal W(u)}=u$ that is analytic on $\C\setminus(-\infty,-1/e]$, real and increasing on $(-1/e,\infty)$, with $\mathcal W(0)=0$, $\mathcal W(-1/e)=-1$ and $|\im\mathcal W|<\pi$. Its Lagrange series $\mathcal W(u)=\sum_{n\ge1}(-n)^{n-1}u^n/n!$ has radius $1/e$, and $\mathcal W(u)=\log u-\log\log u+o(1)$ as $u\to\infty$ \cite[(3.1), (4.19)]{Corless96}. Since $\mathcal W$ is a Bernstein function \cite{Pakes11,KJ10}, $e^{-\mathcal W(u)}=\mathcal W(u)/u$ is the Laplace transform of an infinitely divisible law on $[0,\infty)$, the Lambert law; we write $T_{\mathcal W}$ for a variable with this law. It is a GGC with an explicit, absolutely continuous Thorin measure, and the proof of this fact is a template that fails for the clock at exactly one step.

\begin{proposition}[{\cite[Section 4, (4.4)]{Corless96}}, {\cite[(10), (11)]{KJC11}}]\label{pk:prop:bv}
For $0<b<\pi$ put $r(b)=b\,e^{-b\cot b}/\sin b$. Then $r$ increases from $1/e$ to $\infty$ on $(0,\pi)$, and the boundary values of $\mathcal W$ on the cut, from the upper half plane, are $\mathcal W(-r(b)+i0)=-b\cot b+ib$.
\end{proposition}

\begin{proof}
Put $w=\alpha+ib$ with $0<b<\pi$ in $we^w=-x$, $x>0$. The imaginary part gives $e^\alpha(\alpha\sin b+b\cos b)=0$, so $\alpha=-b\cot b$, and then the real part gives $e^\alpha(\alpha\cos b-b\sin b)=-e^{-b\cot b}\,b/\sin b$, so that $x=r(b)$. Next, $(\log r)'(b)=1/b-2\cot b+b/\sin^2b$ and $b\sin^2b\,(\log r)'(b)=(b-\sin b\cos b)^2+\sin^4b>0$, while $r(0+)=1/e$ and $r(\pi-)=\infty$; so for each $x>1/e$ there is exactly one point $w$ of the curve $\{-b\cot b+ib:0<b<\pi\}$ with $we^w=-x$, namely the point with $r(b)=x$. This curve bounds the range of the principal branch in the upper half plane \cite[Section 4]{Corless96}, and $\mathcal W$ is continuous on the cut from above, with values on it; so $\mathcal W(-r(b)+i0)=-b\cot b+ib$.
\end{proof}

Kalugin, Jeffrey and Corless obtain the next theorem from the Stieltjes representation of $\mathcal W'$ \cite[Theorem 3.3]{KJC11} and the characterisation of Thorin Bernstein functions in \cite[Chapter 8]{boSSV}. The proof below is direct, by boundary values, reflection and a growth bound, and its four steps are compared with the clock at the end of the section.

\begin{theorem}[Kalugin, Jeffrey and Corless {\cite[Theorem 6.2]{KJC11}}, \cite{KJC12}]\label{pk:thm:lambert}
Let $U_{\mathcal W}$ be the image of the uniform law $db/\pi$ on $(0,\pi)$ under $r$, a probability measure on $(1/e,\infty)$. Then for $u\in\C\setminus(-\infty,-1/e]$
\[
\mathcal W(u)=\int\log\Big(1+\frac uz\Big)\,U_{\mathcal W}(dz)=\frac1\pi\int_0^\pi\log\Big(1+\frac u{r(b)}\Big)\,db .
\]
Hence $\mathcal W$ is a Thorin Bernstein function, and $T_{\mathcal W}$ is a GGC with no drift and Thorin measure $U_{\mathcal W}$, of total mass one, whose distribution function is $U_{\mathcal W}((1/e,y])=\pi^{-1}\im\mathcal W(-y+i0)$ for $y>1/e$.
\end{theorem}

\begin{proof}
Let $\Psi(u)$ be the integral. Since $U_{\mathcal W}$ is carried by $(1/e,\infty)$, $\Psi$ is analytic on $\C\setminus(-\infty,-1/e]$, real on $(-1/e,\infty)$, and $\Psi(\bar u)=\overline{\Psi(u)}$; so is $\mathcal W$. Let $x>1/e$, and let $b(x)\in(0,\pi)$ be the point with $r(b(x))=x$. As $\eta\downarrow0$, $1+(-x+i\eta)/z$ tends to a negative number from the upper half plane when $z<x$ and to a positive number when $z>x$, and $U_{\mathcal W}$ has no atoms, so by dominated convergence $\im\Psi(-x+i0)=\pi\,U_{\mathcal W}((1/e,x))=b(x)$; by Proposition~\ref{pk:prop:bv} also $\im\mathcal W(-x+i0)=b(x)$, which is the formula for the distribution function. Near $1/e$, $r(b)-1/e\sim b^2/(2e)$, so $U_{\mathcal W}$ has a density asymptotic to $(e/2)^{1/2}\pi^{-1}(z-1/e)^{-1/2}$ there, and a bounded continuous density on $[1/e+c,\infty)$ for every $c>0$. Hence $\Psi$ extends continuously to the closed upper half plane, and so does $\mathcal W$, with $\mathcal W(-1/e)=-1$. The difference $D=\mathcal W-\Psi$ is therefore continuous on the closed upper half plane and real on the real axis, where on $(-\infty,-1/e)$ the imaginary parts of the two boundary values agree. The Schwarz reflection principle continues $D$ across $(-\infty,-1/e)$, the continuation coincides with $D$ in the lower half plane because $D(\bar u)=\overline{D(u)}$, and the singularity at $-1/e$, where $D$ is bounded, is removable. So $D$ is entire, and $D(0)=0$.

It remains to bound $D$, and one sided bounds on real parts suffice. First, $\re\Psi(u)=\int\log|1+u/z|\,U_{\mathcal W}(dz)\le\log(1+e|u|)$, since $z>1/e$. Next let $|u|=R\ge e$ and $w=\mathcal W(u)$, so that $|w|e^{\re w}=R$ and $|\im w|<\pi$. If $|w|>\log R+\pi$, then $|\re w|>|w|-\pi>\log R$. This is impossible when $\re w\ge0$, since then $\re w=\log R-\log|w|<\log R$; and when $\re w<0$ it gives $\log|w|-\log R>|w|-\pi$, that is $|w|-\log|w|<\pi-\log R\le\pi-1$, which fails for $|w|>1+\pi$, since $x-\log x$ increases for $x>1$ and exceeds $\pi-1$ at $x=1+\pi$. Hence $|w|\le\log R+\pi$ and $\re\mathcal W(u)=\log R-\log|w|\ge\log R-\log(\log R+\pi)$. By continuity both bounds hold on the whole circle, its point on the cut included, so
\[
\max_{|u|=R}\re\big(-D(u)\big)\le\log(1+eR)-\log R+\log(\log R+\pi),\qquad R\ge e .
\]
The Borel Carath\'eodory lemma for the entire function $-D$, which vanishes at the origin, gives $\max_{|u|=R/2}|D(u)|\le2\max_{|u|=R}\re(-D(u))$, which is $O(\log\log R)$ by this bound. By Cauchy's estimates every Taylor coefficient of $D$ of positive order vanishes, so $D\equiv D(0)=0$.
\end{proof}

The growth step cannot be done by writing the integral as $\log u$ plus a bounded term. The total mass of $U_{\mathcal W}$ is one, but $\int\log z\,U_{\mathcal W}(dz)=\infty$, since $U_{\mathcal W}((y,\infty))\sim1/\log y$ as $y\to\infty$, and $\mathcal W(u)-\log u=-\log\log u+o(1)$ is unbounded; the one sided bounds avoid this.

\begin{corollary}[{first part from \cite[Theorem 6.5]{KJC11} and \cite{KJC12}}]\label{pk:cor:levy}
The canonical function of $T_{\mathcal W}$ is
\[
k_{\mathcal W}(x)=\int e^{-xz}\,U_{\mathcal W}(dz)=\frac1\pi\int_0^\pi e^{-x\,r(b)}\,db,
\]
so that $\mathcal W(u)=\int_0^\infty(1-e^{-ux})\,k_{\mathcal W}(x)\,dx/x$ for $u\ge0$. The symmetric law of $Z\sqrt{2T_{\mathcal W}}$, with $Z$ standard normal and independent of $T_{\mathcal W}$, has characteristic function $e^{-\mathcal W(w^2)}$, is an EGGC (Section~\ref{sec:ch3:hierarchy}), and has L\'evy density $k^{\rm s}_{\mathcal W}(|x|)/|x|$ with
\[
k^{\rm s}_{\mathcal W}(x)=\int e^{-x\sqrt z}\,U_{\mathcal W}(dz)=\frac1\pi\int_0^\pi e^{-x\sqrt{r(b)}}\,db,
\]
that is, $\mathcal W(w^2)=\int_0^\infty4\sin^2(wx/2)\,k^{\rm s}_{\mathcal W}(x)\,dx/x$ for real $w$.
\end{corollary}

\begin{proof}
Insert Frullani's integral $\log(1+u/z)=\int_0^\infty(1-e^{-ux})e^{-zx}x^{-1}dx$ in Theorem~\ref{pk:thm:lambert} and exchange the integrals by Tonelli's theorem. For the symmetric law, $\E e^{iwZ\sqrt{2T_{\mathcal W}}}=\E e^{-w^2T_{\mathcal W}}$, and the cosine dictionary \eqref{rz:eq:cosdict} with $U=U_{\mathcal W}$, for which $K_U=2k^{\rm s}_{\mathcal W}$, gives the L\'evy density $K_U(|x|)/(2|x|)$.
\end{proof}

\begin{remark}[Self decomposability, after Pakes]\label{pk:rem:sd}
Since $k_{\mathcal W}$ is nonincreasing, $T_{\mathcal W}$ is self decomposable, and as in Proposition~\ref{sb:prop:bdlp} its driving subordinator $\mathsf Z$ has the L\'evy density $-k_{\mathcal W}'(x)=\int ze^{-xz}\,U_{\mathcal W}(dz)$, of total mass $k_{\mathcal W}(0+)=1$ and with mean $\int_0^\infty k_{\mathcal W}(x)\,dx=\int z^{-1}\,U_{\mathcal W}(dz)=1$ (Corollary~\ref{pk:cor:nuttall}). So the driving process is compound Poisson with unit rate and jumps $J$ of mean one; its Laplace exponent is $u\mathcal W'(u)=\mathcal W(u)/(1+\mathcal W(u))$, and $\E e^{-uJ}=1/(1+\mathcal W(u))$ \cite[(5), (6), (18)]{Pakes26}. Stopping $T_{\mathcal W}=\int_0^\infty e^{-t}\,d\mathsf Z_t$ at the first jump of $\mathsf Z$, at an exponential time $\tau$, gives the perpetuity $T_{\mathcal W}\stackrel d=\mathsf U(T_{\mathcal W}+J)$, with $\mathsf U=e^{-\tau}$ uniform on $(0,1)$ and the three variables independent. The stationary excess $\mathcal SJ$ has the density $P(J>x)$, since $\E J=1$, and the Laplace transform $(1-\E e^{-uJ})/u=\mathcal W'(u)$; since $\E e^{-u\mathsf UY}=u^{-1}\int_0^u\E e^{-sY}ds$ for $Y\ge0$ independent of $\mathsf U$, also $T_{\mathcal W}\stackrel d=\mathsf U\,\mathcal SJ$, and $\mathcal SJ\stackrel d=J+T_{\mathcal W}$. The converse, part (b) of \cite[Theorem 1]{Pakes26}, has the same short proof, which we repeat in our notation. If $X\stackrel d=\mathsf U(X+J)$ and $X\stackrel d=\mathsf U\,\mathcal SJ$ for some $J\ge0$ of mean $\mu\in(0,\infty)$, put $\varphi(u)=\E e^{-uX}$, $j(u)=\E e^{-uJ}$ and $F(u)=u\varphi(u)$. The two relations say that $F'=\varphi j$ and $F'=(1-j)/(\mu u)$; eliminating $j=uF'/F$ gives $F'=F/(u(1+\mu F))$, so $y(u)=\mu F(u/\mu)$ satisfies $y'=y/(u(1+y))$, that is $(\log y+y)'=1/u$, with $y(u)\sim u$ as $u\downarrow0$. Hence $ye^y=u$ and $y=\mathcal W$ on $(0,\infty)$, which gives $\E e^{-uX}=\mathcal W(\mu u)/(\mu u)$, so $X\stackrel d=\mu T_{\mathcal W}$, and $j(u)=1/(1+\mathcal W(\mu u))$, so $J/\mu$ has the law of the jumps above.
\end{remark}

\begin{corollary}[{\cite[Corollary 3.4]{KJC11}}, \cite{KJCB12}; Nuttall's integral \cite{Nuttall85,Bouwkamp86}]\label{pk:cor:nuttall}
For $n\ge0$,
\[
\int z^{-n}\,U_{\mathcal W}(dz)=\frac1\pi\int_0^\pi\Big(\frac{\sin b\;e^{b\cot b}}b\Big)^n\,db=\frac{n^n}{n!}\qquad(0^0=1),
\]
so $(n^n/n!)_{n\ge0}$ is the moment sequence of the image of $U_{\mathcal W}$ under $z\mapsto1/z$, a probability measure on $(0,e)$. Nuttall conjectured, and Bouwkamp proved, the identity with $\nu^\nu/\Gamma(\nu+1)$ for every real $\nu\ge0$.
\end{corollary}

\begin{proof}
Differentiating Theorem~\ref{pk:thm:lambert} gives
\[
\mathcal W'(u)=\int\frac{U_{\mathcal W}(dz)}{u+z}=\sum_{n\ge1}(-u)^{n-1}\int z^{-n}\,U_{\mathcal W}(dz),\qquad |u|<1/e,
\]
the series converging uniformly in $z$ since $z>1/e$. Compare with the derivative of the Lagrange series $\mathcal W(u)=\sum_{n\ge1}(-n)^{n-1}u^n/n!$, whose coefficients are $(-n)^{n-1}/(n-1)!=(-1)^{n-1}n^n/n!$. For $n=0$ the integral is the total mass, one.
\end{proof}

Conversely, as Kalugin, Jeffrey and Corless observe after their Corollary 3.4 \cite{KJC11}, Nuttall's integral, which Bouwkamp proved without the Lambert function, gives back the theorem. It says that the image $\mu$ of $U_{\mathcal W}$ under $z\mapsto1/z$ has the moments $n^n/n!$, so $\int t\,\mu(dt)/(1+ut)=\int U_{\mathcal W}(dz)/(u+z)$ and $\mathcal W'(u)$ have the same Taylor series $\sum_{m\ge0}(-u)^m(m+1)^m/m!$ at the origin; both are analytic on $\C\setminus(-\infty,-1/e]$, so they agree, and integrating from $0$ gives Theorem~\ref{pk:thm:lambert}. Pakes's second proof that $\mathcal W$ is Thorin Bernstein also goes through Nuttall's integral \cite{Pakes18}, \cite[Section 2]{Pakes26}.

The program \texttt{pakes\_lambert.py} checks the section. At $40$ digits the boundary values of Proposition~\ref{pk:prop:bv} hold to $3\times10^{-40}$ at eight values of $b$ from $0.05$ to $3.1$ (at $30$ digits, at $b=1.3$, to all digits), the identity for $b\sin^2b\,(\log r)'$ holds to $10^{-39}$, and $(r(b)-1/e)/b^2$ is $0.18393972059$ at $b=10^{-6}$, against $1/(2e)$. The integral of Theorem~\ref{pk:thm:lambert} agrees with $\mathcal W$ to $2\times10^{-40}$ at seventeen points, ten real ones from $-0.3678$ to $10^{12}$ and seven complex ones, some next to the cut, such as $-5+0.001i$ and $-10^6+10^3i$. At $60$ digits Newton's method on $we^w=27/10$, the integral and Arb's Lambert function agree: $\mathcal W(2.7)=0.99662873426547746603029759804008309\ldots$. The binary double nearest $2.7$ exceeds $2.7$ by $1.8\times10^{-16}$, and $\mathcal W$ there is larger by $3.3\times10^{-17}$, so for digits beyond the sixteenth the input must be passed exactly, as the rational $27/10$. The moments of Corollary~\ref{pk:cor:nuttall} agree with $n^n/n!$ for $n\le12$ to $10^{-37}$, and with $\nu^\nu/\Gamma(\nu+1)$ at $\nu=0.5$, $1.5$, $2.7$ and $7.25$ to $3\times10^{-39}$; for $n=0,\dots,6$ they are $1$, $1$, $2$, $4.5$, $10.6667$, $26.0417$ and $64.8$. In double precision the L\'evy integrals of Corollary~\ref{pk:cor:levy} reproduce $\mathcal W(u)$ at $u=1$, $2.7$ and $10$, and $\mathcal W(w^2)$ at $w=1$ and $2$, to $6\times10^{-16}$, and $k_{\mathcal W}(x)$ is $0.988$ and $0.9988$ times $\sqrt{e/2}\,e^{-x/e}/\sqrt{\pi x}$ at $x=100$ and $1000$, as the density of $U_{\mathcal W}$ at $1/e$ predicts; at $z-1/e=1.8\times10^{-9}$ that density is $0.999999995$ times its asymptote. The growth facts behind the proof are visible too: $\mathcal W(u)-\log u+\log\log u$ falls from $0.308$ at $u=10^2$ to $0.0194$ at $10^{128}$; $U_{\mathcal W}((y,\infty))$ is $0.0290$ at $y=10^{16}$, against $1/\log y=0.0271$; and the part of $\int\log z\,U_{\mathcal W}(dz)$ from $b<\pi-c$ exceeds $\log(1/c)$ by $0.126$, $0.144$, $0.145$ and $0.145$ at $c=10^{-2}$, $10^{-4}$, $10^{-8}$ and $10^{-16}$. On the circles $|u|=R$ the lower bound $\re\mathcal W(u)\ge\log R-\log(\log R+\pi)$ of the proof holds with the margins $0.81$, $0.76$, $0.65$, $0.47$ and $0.30$ at $R=3$, $10$, $10^2$, $10^4$ and $10^8$, attained on the cut, and $|\mathcal W(u)|\le0.74(\log R+\pi)$ there.

What transfers to the clock. The proof of Theorem~\ref{pk:thm:lambert} has four steps: the boundary values on the cut, the positivity of the boundary measure, the identification by reflection and a growth bound, and, for the symmetric law, the cosine dictionary \eqref{rz:eq:cosdict}. For the clock the first step is the argument count. Continued along the upper side of the negative axis, past the zeros of $G$ on it by small half circles in the upper half plane, $\im\log G(-x+i0)=\pi N_{\rm line}(\sqrt x)$, where $N_{\rm line}(t)$ counts with multiplicity the zeros on the critical line with ordinate in $(0,t)$, because the four zeros $\rho$, $1-\rho$, $\bar\rho$, $1-\bar\rho$ of a zero off the line contribute to $G$ the factor $(1+u/a)(1+u/\bar a)=|1+u/a|^2>0$ at real $u$; so $N_{\rm line}(\sqrt x)$, the distribution function of $U_{\rm line}$ (Proposition~\ref{prop:struct}), is read from the boundary values of $\log G$ with no hypothesis, and it is $N(\sqrt x)$ under RH. The second step is then automatic, since a counting function is nondecreasing. The fourth step gives the Poisson trace $K_{1/2}(y)=\sum_ke^{-y\sqrt{a_k}}$, which is $\sum_{\gamma>0}e^{-\gamma y}$ under RH, as the L\'evy side of $Y_{1/2}=Z\sqrt{2T}$ (Section~\ref{sm:sec:poisson}), while the heat trace $W$ is the canonical function of $T$ itself (Section~\ref{sec:ch3:readings}); the two traces are linked with no hypothesis by subordination. The third step is the obstruction: $\mathcal W$ is analytic off $(-\infty,-1/e]$ by construction, whereas $\log G$ is analytic off $(-\infty,0]$ if and only if RH holds (Proposition~\ref{prop:rigid} and the argument after it). In the Lambert case the boundary measure is a continuum lying on the cut; for the clock its atoms are the zeros, and whether they all lie on the cut is the whole question.
